% BuzzCalculus 技術報告 —— 用 WSL 的 xelatex 編（跑兩次才有目錄）：
%   wsl -e bash -lc "cd /mnt/c/Users/Tudo/Downloads/BuzzCalculus/docs/report && xelatex -interaction=nonstopmode buzzcalculus_engines.tex && xelatex -interaction=nonstopmode buzzcalculus_engines.tex"
% 所有數字取自 2026-09-16 的程式碼與驗證器輸出（見各章註記）。
\documentclass[11pt,a4paper,oneside]{report}

\usepackage[a4paper,top=2.4cm,bottom=2.6cm,left=2.4cm,right=2.4cm]{geometry}
\usepackage{fontspec}
\usepackage{xeCJK}
\usepackage{xcolor}
\usepackage{graphicx}
\usepackage{booktabs}
\usepackage{longtable}
\usepackage{tabularx}
\usepackage{multirow}
\usepackage{array}
\usepackage{enumitem}
\usepackage{amsmath,amssymb}
\usepackage{tikz}
\usetikzlibrary{arrows.meta,positioning,calc}
\usepackage[most]{tcolorbox}
\usepackage{listings}
\usepackage{caption}
\usepackage{titlesec}
\usepackage{fancyhdr}
\usepackage{microtype}
\usepackage[hidelinks]{hyperref}

% ── 字型：拉丁字用 TeX Gyre，中文用 Noto CJK TC，等寬用 DejaVu ──
\setmainfont{TeX Gyre Pagella}
\setsansfont{TeX Gyre Heros}
\setmonofont{DejaVu Sans Mono}[Scale=0.86]
\setCJKmainfont{Noto Serif CJK TC}[BoldFont={Noto Serif CJK TC Bold}]
\setCJKsansfont{Noto Sans CJK TC}
\setCJKmonofont{Noto Sans CJK TC}
\xeCJKsetup{CJKmath=true,PunctStyle=kaiming}
\linespread{1.32}

% ── 色彩：紙、墨、金、苔綠、鐵鏽 ──
\definecolor{ink}{HTML}{1F1D18}
\definecolor{muted}{HTML}{625D52}
\definecolor{gold}{HTML}{B98A22}
\definecolor{goldsoft}{HTML}{F5EBD0}
\definecolor{moss}{HTML}{2F6B4A}
\definecolor{mosssoft}{HTML}{E4EEE6}
\definecolor{rust}{HTML}{A8432C}
\definecolor{rustsoft}{HTML}{F5E0D8}
\definecolor{paper}{HTML}{F6F2E8}
\definecolor{line}{HTML}{D9D2C2}
\color{ink}

% ── 章節標題 ──
\titleformat{\chapter}[display]
  {\normalfont\sffamily}
  {\color{gold}\small\bfseries\chapterlabel}
  {6pt}
  {\color{ink}\Huge\bfseries}
  [\vspace{4pt}{\color{gold}\titlerule[2pt]}]
\titlespacing*{\chapter}{0pt}{-10pt}{26pt}
\newcommand{\chapterlabel}{第\,\thechapter\,章}
\renewcommand{\chaptermark}[1]{\markboth{\chapterlabel\quad #1}{}}
\titleformat{\section}{\normalfont\sffamily\Large\bfseries\color{ink}}{\thesection}{0.8em}{}
\titleformat{\subsection}{\normalfont\sffamily\large\bfseries\color{ink}}{\thesubsection}{0.8em}{}
\titlespacing*{\section}{0pt}{18pt}{8pt}
\titlespacing*{\subsection}{0pt}{12pt}{5pt}

% ── 頁首頁尾 ──
\pagestyle{fancy}
\fancyhf{}
\fancyhead[L]{\small\sffamily\color{muted}BuzzCalculus 技術報告}
\fancyhead[R]{\small\sffamily\color{muted}\nouppercase{\leftmark}}
\fancyfoot[C]{\small\sffamily\color{muted}\thepage}
\renewcommand{\headrulewidth}{0.4pt}
\renewcommand{\headrule}{\vspace{-2pt}{\color{line}\hrule width\headwidth height 0.4pt}}
\fancypagestyle{plain}{\fancyhf{}\fancyfoot[C]{\small\sffamily\color{muted}\thepage}\renewcommand{\headrulewidth}{0pt}}

% ── 呼叫框：設計決定、坑、數字 ──
\newtcolorbox{decision}[1][設計決定]{enhanced,breakable,colback=goldsoft,colframe=gold,boxrule=0.6pt,arc=2pt,
  title=#1,fonttitle=\sffamily\bfseries\small,coltitle=ink,attach boxed title to top left={yshift=-2mm,xshift=4mm},
  boxed title style={colback=goldsoft,colframe=gold,boxrule=0.6pt,arc=1pt},top=6pt,left=8pt,right=8pt,bottom=6pt}
\newtcolorbox{pitfall}[1][實際踩過的坑]{enhanced,breakable,colback=rustsoft,colframe=rust,boxrule=0.6pt,arc=2pt,
  title=#1,fonttitle=\sffamily\bfseries\small,coltitle=ink,attach boxed title to top left={yshift=-2mm,xshift=4mm},
  boxed title style={colback=rustsoft,colframe=rust,boxrule=0.6pt,arc=1pt},top=6pt,left=8pt,right=8pt,bottom=6pt}
\newtcolorbox{principle}[1][原則]{enhanced,breakable,colback=mosssoft,colframe=moss,boxrule=0.6pt,arc=2pt,
  title=#1,fonttitle=\sffamily\bfseries\small,coltitle=ink,attach boxed title to top left={yshift=-2mm,xshift=4mm},
  boxed title style={colback=mosssoft,colframe=moss,boxrule=0.6pt,arc=1pt},top=6pt,left=8pt,right=8pt,bottom=6pt}

% ── 程式碼 ──
\lstdefinelanguage{JavaScript}{keywords={const,let,var,function,return,of,for,if,else,true,false,null,new,await,async,while,continue,break},sensitive=true,comment=[l]{//},string=[b]"}
\lstset{
  basicstyle=\ttfamily\footnotesize\color{ink},
  keywordstyle=\color{moss}\bfseries,
  commentstyle=\color{muted},
  stringstyle=\color{rust},
  backgroundcolor=\color{paper},
  frame=single, framerule=0.3pt, rulecolor=\color{line},
  xleftmargin=6pt, xrightmargin=6pt, aboveskip=8pt, belowskip=10pt,
  breaklines=true, columns=fullflexible, keepspaces=true, showstringspaces=false,
  language=JavaScript
}
\renewcommand{\contentsname}{目錄}
\renewcommand{\figurename}{圖}
\renewcommand{\tablename}{表}
\renewcommand{\appendixname}{附錄}
\captionsetup{font={small,sf},labelfont={bf,color=gold},skip=6pt}
\setlist{nosep,leftmargin=1.4em}
\newcommand{\code}[1]{\texttt{#1}}
\newcommand{\file}[1]{\texttt{\small #1}}
\newcolumntype{Y}{>{\raggedright\arraybackslash}X}
\newcolumntype{R}{>{\raggedleft\arraybackslash}p{1.6cm}}

\begin{document}

% ════════════════════════ 封面 ════════════════════════
\begin{titlepage}
  \thispagestyle{empty}
  \vspace*{28mm}
  {\color{gold}\sffamily\bfseries\small 技術報告 · 2026 年 9 月\par}
  \vspace{10mm}
  {\sffamily\bfseries\fontsize{34}{42}\selectfont BuzzCalculus 的引擎\par}
  \vspace{6mm}
  {\fontsize{16}{24}\selectfont\color{muted} 沒有後端，怎麼判答案、驗證明、決定你今天練什麼\par}
  \vspace{14mm}
  {\color{gold}\rule{0.32\textwidth}{2pt}\par}
  \vspace{14mm}
  \begin{minipage}{0.9\textwidth}
    \large
    一個純靜態的微積分訓練站：沒有伺服器、沒有帳號，所有運算在瀏覽器裡完成。本報告整理它的三個純函式引擎——答案判分、白話證明檢查、能力模型與抽題——以及把每一條承諾綁進持續整合的驗證文化。全部數字取自 2026 年 9 月 16 日的程式碼與驗證器輸出。
  \end{minipage}
  \vfill
  {\small
  \begin{tabular}{@{}lp{0.78\textwidth}@{}}
    \sffamily\color{muted}專案 & tudohuang / BuzzCalculus \\
    \sffamily\color{muted}線上 & \url{https://tudohuang.github.io/BuzzCalculus/} \\
    \sffamily\color{muted}原始碼 & \url{https://github.com/tudohuang/BuzzCalculus} \\
    \sffamily\color{muted}規模 & 1,994 題 · kernel 14,489 行 · app.js 15,577 行 · 41 支驗證器 · 9 支 E2E \\
  \end{tabular}}
\end{titlepage}

% ════════════════════════ 摘要 ════════════════════════
\chapter*{摘要}
\addcontentsline{toc}{chapter}{摘要}
\markboth{摘要}{}
BuzzCalculus 把「練微積分」這件事拆成三個可以被驗證的問題：使用者打進來的答案對不對、使用者寫的證明有沒有跳步、今天應該練什麼。三個問題各有一個引擎，全部是純函式：輸入紀錄、題庫與時間，輸出資料，不碰 DOM、不碰儲存、不呼叫系統時鐘。這讓同一份程式碼能在瀏覽器裡跑，也能在 Node 的驗證器裡被逐題釘住。

判分引擎不做符號化簡，而是在刻意挑過的取樣點上比較兩個式子的值；不定積分比較的是相對基準點的增量，因此對積分常數不敏感。標準答案本身另有一支獨立驗算器把關：它從題幹推出一條與解題無關的數值路徑再和答案比，目前覆蓋 1,753 題（88\%），不符數在 CI 裡必須為零。

白話證明引擎把中文證明逐行分成 13 種句型，對代數鏈在當前假設下隨機取點驗證，對定理引用比對規則字典；驗不了的三分之一（存在句、抽象函數）標黃而不假裝。每一題的參考證明都經過三種變異測試：刪一行、改一個符號、換方向，變異版必須紅在正確的那一行。

能力模型從作答紀錄回放出 80 個技巧的熟練度，每個常數都能解釋：30 天的權重時間常數、相當於 6 次不確定作答的先驗、加權量 12 才顯示分數。規劃器輸出的不是題目而是「配方＋為什麼」，抽題器抽不滿一定降級補滿並記錄降級路徑。

報告最後一部分是驗證文化：41 支驗證器、9 支以原生 Chrome DevTools Protocol 寫成的端對端測試、隱私政策的靜態檢查、公開數字的自動核對，以及讓 CI 失敗能被看見的 annotation 機制。這些不是任何單一演算法，卻是一個零後端靜態站能被信任的原因。

\tableofcontents

% ════════════════════════ 第 1 章 ════════════════════════
\chapter{導論}
\section{要解的問題}
微積分的練習題網站很多，它們的共同做法是：題目存在伺服器、答案交給伺服器判、進度存在帳號裡。BuzzCalculus 從第一天就拿掉這三件事——沒有伺服器、沒有帳號、沒有任何一筆資料離開使用者的裝置。這不是省錢的妥協，而是產品定位：一個能在考前一小時打開就練、練完關掉不留痕跡、離線也能用的工具。

拿掉後端之後，三個原本理所當然的事變成問題：
\begin{enumerate}[label=\arabic*.]
  \item 使用者打進來的 \code{2*x\string^3}、\code{x\string^3*2}、\code{(2x\string^3)} 都得判對，而題庫存的是 \code{2x\string^3}。判分器必須在瀏覽器裡跑，而且不能誤判——判錯一題造成的信任損失，大於新增五十題的收益。
  \item 「今天練什麼」需要一個能力模型，而能力模型需要作答歷史。沒有資料庫、只有 localStorage 的 5 MB，歷史要怎麼存、怎麼推導、怎麼在使用者中斷幾週之後不「失憶」。
  \item 證明題沒有選項可選、沒有數字可比。瀏覽器裡跑不了 Lean，但「這一步有沒有跳」這件事，真的做不到嗎？
\end{enumerate}

\section{三條設計原則}
\begin{principle}[純函式 kernel]
\file{src/kernel/} 底下 23 個模組、14,489 行，每一個都只吃 \code{(records, problems, now)} 這類輸入、吐資料。不碰 DOM、不碰 localStorage、不呼叫 \code{Date.now()}——時間是參數。這帶來兩個直接的好處：能力曲線可以「回放」到任何一天，所以不需要每天寫快照；同一份程式碼可以被 Node 的驗證器 require 進來逐題釘住。
\end{principle}
\begin{principle}[每一條承諾都有測試]
隱私政策上寫「關掉分析會停止收集」，半年後有人重構把那個判斷拿掉，政策就變成不實陳述，而且沒有任何測試會紅。所以隱私政策裡的每一條、README 裡的每一個數字、每一題的答案、規格文件裡的現況快照，都有對應的靜態檢查：改壞了 CI 就紅。
\end{principle}
\begin{principle}[推薦必須說得出為什麼]
使用者信任推薦系統的唯一來源，是它說得出理由。規劃器輸出的每一份配方都帶一句 \code{why}，它是產品的一部分而不是 debug 訊息，驗證器會擋住空字串。同樣的精神：抽不滿題目時要降級補滿並誠實說明，樣本不夠時能力數字寧可留白也不猜。
\end{principle}

\section{本報告的範圍}
第 2 章講架構與資料流；第 3 章講題庫與內容管線；第 4 到第 8 章是三個引擎（判分與獨立驗算、白話證明、能力模型、規劃與抽題）；第 9 章入門課程；第 10 章手寫計算紙；第 11 章資料層與隱私；第 12 章驗證文化與 CI；第 13 章效能預算與模組化；第 14 章結語。附錄列出所有常數與驗證器。

% ════════════════════════ 第 2 章 ════════════════════════
\chapter{系統架構}
\section{分層}
整個站是一份 \file{index.html}、一份 \file{styles.css}、一支 service worker，和 65 個 \code{<script defer>}。沒有建置步驟，沒有 node\_modules：TypeScript 只在開發與 CI 用 JSDoc 做型別檢查（\code{tsc --noEmit}），不進 runtime。

\begin{table}[h]
\centering\small\sffamily
\begin{tabularx}{\textwidth}{@{}lRY@{}}
\toprule
層 & 行數 & 內容 \\
\midrule
\file{src/app.js} & 15,577 & 畫面、狀態、事件委派。單一 IIFE，正在逐段拆出去（第 13 章）。 \\
\file{src/kernel/*.js}（23 檔） & 14,489 & 純函式引擎與產生的側表：判分取樣、白話證明、能力模型、規劃器、抽題器、紀錄分層、手寫繪製、技巧圖譜、難度三軸、來源、永久編號、驗算側表。 \\
\file{src/problem*.js}（37 檔） & — & 題庫資料，995 KB。 \\
\file{src/course.js}、\file{share\_cards.js}、\file{proof\_lab\_ui.js} & — & 從 app.js 搬出去的整段 render：入門課程、分享卡、證明訓練畫面。 \\
\file{tools/}（66 支） & — & 41 支驗證器、9 支 E2E、產生器、剖析工具。零依賴，只用 Node 內建模組。 \\
\bottomrule
\end{tabularx}
\caption{程式的四層。kernel 與 app 的行數相當，這是刻意的：邏輯要能離開畫面被測。}
\end{table}

\section{一次作答的資料流}
\begin{figure}[h]
\centering
\begin{tikzpicture}[
  node distance=6mm and 5mm,
  box/.style={draw=line,fill=paper,rounded corners=2pt,minimum height=11mm,text width=27mm,align=center,font=\sffamily\small},
  hot/.style={box,draw=gold,line width=1pt},
  arr/.style={-{Latex[length=2mm]},draw=muted,line width=0.8pt}
]
\node[box] (u) {使用者作答\\{\scriptsize\color{muted}式子、選項、手寫}};
\node[hot,right=of u] (c) {判分引擎\\{\scriptsize\color{muted}answer\_sampling}};
\node[hot,right=of c] (a) {能力模型\\{\scriptsize\color{muted}ability}};
\node[hot,right=of a] (p) {規劃器\\{\scriptsize\color{muted}planner}};
\node[box,right=of p] (s) {抽題器\\{\scriptsize\color{muted}session}};
\draw[arr] (u) -- (c); \draw[arr] (c) -- (a); \draw[arr] (a) -- (p); \draw[arr] (p) -- (s);
\draw[arr] (s.south) -- ++(0,-6mm) -| (u.south);
\node[below=9mm of a,font=\scriptsize\color{muted},text width=120mm,align=center]{迴圈只經過 localStorage 與 IndexedDB；能力曲線由紀錄回放推導，不存快照};
\end{tikzpicture}
\caption{每一局訓練的迴圈。三個金框是本報告的主角。}
\end{figure}

\begin{enumerate}[label=\arabic*.]
  \item 使用者送出答案，判分引擎回傳對錯與一句給人看的訊息（不是判分器的內部話）。
  \item 作答寫進三層紀錄（第 11 章）：完整明細、精簡逐題 tuple、局摘要。
  \item 能力模型從紀錄回放出每個技巧的熟練度、置信度、速度象限。
  \item 規劃器據此產生配方：幾題到期複習、幾題弱點、幾題新技巧，以及為什麼。
  \item 抽題器照配方抽，抽不滿就降級補滿，並把降級寫進 \code{meta.fallbacks}。
\end{enumerate}

\section{離線與載入}
service worker 精確列出快取清單（101 個項目涵蓋 65 支 script），\file{validate\_offline\_assets.js} 檢查 app 用到的每一個圖示都在自帶的圖示集裡、每一個資產都在快取清單裡；\file{validate\_app\_shell.js} 檢查 repo 裡每一個要出貨的 HTML 都被 CI 的 Assemble 步驟複製了。CACHE\_NAME 必須以 app 版本號開頭、後面接日期或標籤，否則同一版內改資源使用者拿不到新檔（\file{validate\_version.js}）。

% ════════════════════════ 第 3 章 ════════════════════════
\chapter{題庫與內容管線}
\section{規模與分佈}
題庫 1,994 題，只有微積分（理科題已於 2026 年 8 月移到姊妹站）。四個主題、六個難度等級：

\begin{table}[h]
\centering\small\sffamily
\begin{tabular}{@{}lrrrrrrr@{}}
\toprule
主題 & R1 & R2 & R3 & R4 & R5 & R6 & 合計 \\
\midrule
極限 & 45 & 60 & 51 & 40 & 66 & 69 & 331 \\
微分 & 53 & 116 & 132 & 167 & 111 & 110 & 689 \\
積分 & 39 & 77 & 85 & 129 & 131 & 243 & 704 \\
級數 & 33 & 42 & 45 & 25 & 60 & 65 & 270 \\
\midrule
合計 & 170 & 295 & 313 & 361 & 368 & 487 & 1,994 \\
\bottomrule
\end{tabular}
\caption{題庫依主題與難度的分佈。R6 最多，因為競賽風格與名校風格的題包集中在那裡。}
\end{table}

答案型別決定判分路徑：數值 1,356 題、式子 317、不定積分 142、文字判定 100、集合 17、區間 16、作圖表 22、互動圖形 24（點位 8、切線 4、選圖 12）。

\section{難度不是一個數字，是三軸}
每一題的難度由三個 1–3 的軸決定：步數（Steps）、冷門度（Obscurity）、負荷（Load，題幹與答案的 token 數加上是否需要動筆）。等級由三軸算出，公式在 \file{kernel/rubric.js}，瀏覽器與工具鏈共用同一份——曾經有過兩份，結果畫面上的 R5 和難度說明頁算出來的 R3 對不起來。

\begin{lstlisting}
function rankFrom(axes, skillCount) {
  const raw = axes[0] + axes[1] + axes[2];          // 3..9
  let rank = RAW_TO_RANK[raw] || (raw <= 3 ? 1 : 6);
  // Obscurity at 3 with a high total: at least R5. Only this axis counts --
  // many steps means "takes time", not "cannot do".
  if (axes[1] === 3 && raw >= 7) rank = Math.max(rank, 5);
  // Chaining three or more independent skills: +1. Each step is easy;
  // the difficulty is thinking to connect them.
  if (skillCount >= 3) rank = Math.min(6, rank + 1);
  return rank;
}
\end{lstlisting}

\begin{decision}
三軸表是從現行等級反推的，遷移當下每一題的等級完全不變。表裡標成人工複核過的題代表有人真的看過並負責；其餘的理由由程式依三軸生成，而且會明講「未經人工複核」。機器推導的東西不能看起來像有人背書。
\end{decision}

\section{來源聲明與永久編號}
每一題有一筆來源聲明：原創、改編、取材自某種風格、公有領域。\file{validate\_content\_metadata.js} 掃題幹、解說與來源欄位裡的名校與競賽字眼（大小寫敏感、字界比對——否則 \code{limit} 裡的 \code{mit} 會讓整個題庫誤報），出現的題不准標成原創，而且來源註只准寫「某某風格的原創題」。2026 年 9 月起產品可見文字裡連競賽名稱都不出現。

每一題另有一個永不回收的永久編號（\file{kernel/uid\_map.js}）。它存在的理由是對戰碼：碼裡裝的是題目的永久編號而不是隨機種子，題庫成長重排之後，兩台裝置解出來仍是同一份題。

\section{等價題與模板}
不同題包裡常有只差記法的同一題（\code{\textbackslash ln} 與 \code{\textbackslash log}、\code{x\string^2/2} 與 \code{\textbackslash frac\{x\string^2\}\{2\}}）。刪任何一邊都會在那個題包上開一個洞，所以改成抽題時去重：\file{detect\_duplicates.js} 產生等價側表（42 組，其中 3 組人工複核判定為誤判），同一局或同一份考卷只會出現其中一題。另有模板題（\code{tmpl-*}）由參數展開，\file{expand\_templates.js --check} 保證展開結果與版控一致。

% ════════════════════════ 第 4 章 ════════════════════════
\chapter{引擎一：答案判分}
\section{為什麼不做符號化簡}
判斷兩個式子等價的正統做法是符號化簡：把兩邊都化成標準形再比。瀏覽器裡沒有電腦代數系統，自己寫一個又會在三角恆等式、對數律、絕對值上無窮無盡地補洞。BuzzCalculus 選了另一條路：在一組固定的取樣點上求值，逐點比較。這是有限取樣，不是形式證明——程式碼第一行就寫著這句話——但它的失敗方式是可預測的，而且可以用取樣點的選擇把誤判壓到實務上看不到。

\section{取樣點的選擇}
\begin{lstlisting}
// Positive-only samples had a real hole: sqrt(x^2) and x agree on every
// positive point but are not the same function. Negatives separate them.
// "Nice" numbers (0.5, 1, 2) are where different functions collide most
// often (sin(pi/6) = 1/2 and friends), so the points sit near irrationals.
const SAMPLES = [0.3137, 0.7211, 1.2345, 1.9871, 3.3013,
                 -0.6180, -1.3247, -2.1069];
\end{lstlisting}
單變數用這 8 個點；多變數用線性同餘產生 24 個點，每個座標獨立產生，避免全落在 $y = x + c$ 一條直線上。題目可以帶定義域約束（例如 $x > 0$），取樣點先過濾；至少要有三個點兩邊都有定義才算「判得穩」，否則回「格式讀不穩」而不是判錯。

\section{容差}
\[
  |a - b| \le \max\bigl(10^{-6},\ 10^{-5}\,|a|\bigr)
\]
相對容差 $10^{-5}$ 對手算的式子綽綽有餘，同時足以分開 $\sin x$ 與它的三階 Taylor 多項式在 $x = 1.9871$ 的差。

\section{不定積分：比增量，不比值}
反導函數差一個常數也算對，所以不能直接比 $F(x)$。取第一個取樣點當基準，比較每一點相對基準的增量：
\[
  \bigl|(b_i - b_0) - (a_i - a_0)\bigr| \le \max\bigl(10^{-5},\ 10^{-5}\,|a_i - a_0|\bigr) + 4\varepsilon_{\text{mach}}\max(|b_i|, |b_0|)
\]
容差由標準答案的變化量決定，不能被使用者寫進去的巨大常數放大；最後一項是浮點捨入的保護——常數大到捨入誤差超過容差時，明講「積分常數太大，無法可靠比較」。

\begin{table}[h]
\centering\small\sffamily
\begin{tabularx}{\textwidth}{@{}lYl@{}}
\toprule
題目 & 學生輸入 & 判定 \\
\midrule
$\int 6x^2\,dx$ & \code{2x\string^3 + 100} & \textcolor{moss}{\textbf{正確}}：原函數相差常數 \\
$\int 6x^2\,dx$ & \code{x\string^3} & \textcolor{rust}{\textbf{錯}}：微分後不相同，附「差一個 1/3」的提示 \\
$\frac{d}{dx}\sqrt{x}$ & \code{1/(2*sqrt(x))} 與 \code{0.5*x\string^(-0.5)} & \textcolor{moss}{\textbf{正確}}：多點代入等價 \\
$\lim_{x\to 0}\frac{\sin x}{x}$ & \code{1}、\code{1.0}、\code{2/2} & \textcolor{moss}{\textbf{正確}} \\
\bottomrule
\end{tabularx}
\caption{判分引擎的實際行為。訊息一律是給人看的話，不是判分器的內部話。}
\end{table}

\section{其他答案型別}
\begin{itemize}
  \item \textbf{集合}：元素用逗號隔開，順序不影響判分，每個元素各自走數值比較。
  \item \textbf{區間}：端點用 \code{(} \code{)} 或 \code{[} \code{]}，開閉有差；聯集用 \code{U}；無窮寫 \code{inf}。格式說明在輸入框旁邊先講清楚，不讓使用者為了猜格式重打。
  \item \textbf{文字判定}：收斂、發散、條件收斂、絕對收斂、不存在，含中英別名。
  \item \textbf{互動圖形}：點位題比較點到目標的距離，切線題比較拖出的斜率與正解，都有容差且兩端都判。
  \item \textbf{作圖表}：每格一個判分，錯的選項都要具名說明。
\end{itemize}

\section{選擇題的誘答理由}
選擇題答錯時，「在 $x = 0.31$ 代入時不相同」是判分器的話，不是給人看的。實作改成：告訴使用者選了哪一個、它像是怎麼錯的（誘答理由，例如「指數減 1 忘了」），正解在下面。誘答選項在產生時已經用判分器濾掉與正解等價的寫法，否則會出現兩個都對的選項。

\section{輸入語法}
答案輸入採 WebWork 式語法（\code{*} 乘、\code{\string^} 次方、\code{sqrt}、\code{log} 即 $\ln$）。輸入框有即時 LaTeX 預覽，回答的是「我打的被讀成什麼」——那是送出前唯一能自我檢查的地方。手機上的數學鍵盤依題型排序：積分題先看到 \code{\string^ / ( ) sqrt}，極限題先看到 \code{pi e inf}。

% ════════════════════════ 第 5 章 ════════════════════════
\chapter{標準答案本身的獨立驗算}
判分引擎只能保證「跟標準答案一致」。標準答案錯了怎麼辦？在 \file{verify\_answers.js} 之前，答案的正確性完全靠人工檢查——也就是完全沒有機器把關。

\section{做法}
從\textbf{題幹}推出一條與解題無關的數值路徑，再跟標準答案比：
\begin{itemize}
  \item 導數：對題幹的函數做數值微分，與答案的式子在取樣點比。
  \item 不定積分：把答案微分回去，與被積函數比。
  \item 定積分、極限、級數：直接用數值方法算出來，與答案的值比。
  \item 專用驗算：收斂半徑、Taylor 係數、Hessian 行列式、多變數極限、留數、Wronskian、區域積分、參數曲線、卷積、Beta 與 Gamma 積分、常微分方程的殘差……共 90 餘種方法。
\end{itemize}

\begin{table}[h]
\centering\small\sffamily
\begin{tabular}{@{}lr@{\hspace{14mm}}lr@{}}
\toprule
方法 & 題數 & 方法 & 題數 \\
\midrule
定積分 & 303 & 收斂半徑 & 41 \\
極限 & 257 & 級數收斂判定 & 42 \\
導數（式子） & 202 & 參數積分 & 36 \\
不定積分 & 142 & Taylor 係數 & 33 \\
級數和 & 131 & 線積分 & 26 \\
Taylor 導數 & 57 & 區域積分 & 18 \\
數值微分（點值） & 32 & 一階常微分方程 & 16 \\
\bottomrule
\end{tabular}
\caption{驗算方法的分佈（前 14 種）。全部 1,753 題通過，241 題目前無法獨立驗算（多為文字判定與概念題）。}
\end{table}

\section{CI 門檻}
\code{--ci} 模式下兩件事會讓建置失敗：不符數大於零，或覆蓋率低於門檻。驗算結果寫成側表 \file{kernel/verified\_answers.js}，題庫頁上每一題有「答案已驗算」的標記——這是競品沒有的一句話，藏在 CI 裡等於沒有。

\begin{pitfall}[危險的是安靜地算錯]
驗算器最危險的失敗方式不是「算不出來」，是「安靜地算錯」：一支方法對某類題幹解析錯了但仍回傳一致，覆蓋率好看、其實沒驗。所以黃金測試（\file{validate\_golden.js}）錄下每一支驗算方法在固定題目上的中間結果，改動解析器就會紅。
\end{pitfall}

% ════════════════════════ 第 6 章 ════════════════════════
\chapter{引擎二：白話證明檢查}
\section{三分之二的誠實}
瀏覽器裡跑不了 Lean，Lean 只在 CI 裡驗 8 題經典證明。但「這一步有沒有跳」這件事，做得到三分之二——前提是老實分清楚哪三分之二：
\begin{enumerate}[label=\arabic*.]
  \item \textbf{句型引擎}：認得這句話在做什麼（設變數、假設、代數推導、引用定理、分情況、歸納、下結論）。
  \item \textbf{代數引擎}：$A = B \le C < D$ 這種鏈，在目前的假設下隨機取點驗——等式看差、不等式看方向。
  \item \textbf{推論引擎}：一小本規則字典（均值定理、三角不等式、夾擠……），比對「你寫的句子」與「規則會吐出的形狀」。
\end{enumerate}
剩下的三分之一——存在句、抽象函數 $f$、數論——標黃：看得懂但驗不了，不假裝。

\section{正規化}
使用者會混用全形標點、希臘字母、LaTeX 指令與 Unicode 符號。第一步把它們收斂成一種寫法：全形數字與字母轉半形；「，。：；（）」轉半形；$\le$、$\ge$、$\ne$、$-$、$\times$、$\cdot$ 轉成 ASCII；$\varepsilon$ 與 \code{\textbackslash epsilon} 都變 \code{eps}；$\sqrt{\ }$、$\pi$、$\infty$ 各有對應；\code{|x-2|} 這種不巢狀的絕對值翻成 \code{abs(x-2)}。

\section{十三種句型}
\begin{table}[h]
\centering\small\sffamily
\begin{tabularx}{\textwidth}{@{}llY@{}}
\toprule
句型 & 標籤 & 觸發的寫法（節錄） \\
\midrule
let & 任取／設／取 & 任取、任意取、給定、固定、設、令、取、let、fix、choose \\
assume & 假設／若…則… & 假設、若、如果、suppose、assume、if \\
by & 由 <規則>，… & 由、根據、依據、利用、by、using、applying \\
because & 因為…，所以… & 因為…，所以／故／因此 \\
claim & 則／所以 <推導> & 則、那麼、得、所以、故、因此、於是、即、同理 \\
cases / case & 分情況、情況 N & 分兩種情況、考慮三種情況、case 1 \\
induction & 用歸納法 & 用數學歸納法、對 n 做歸納 \\
base / hypothesis & 當 n = 1 時、假設 n = k 時成立 & 當 n = 1 時、假設 n = k 時 \\
contradiction-start / contradiction & 反設、矛盾 & 反設、假設不然、suppose not；…矛盾 \\
qed & 得證 & 得證、證畢、Q.E.D.、$\blacksquare$ \\
unknown & — & 一行有兩句、或只有式子沒有句型 \\
\bottomrule
\end{tabularx}
\caption{句型引擎的分類。一行只能一句：一行裡有兩個句號就直接判成讀不懂，並說為什麼。}
\end{table}

\section{代數鏈的驗證}
一行 \code{則 |3x − 6| = 3|x − 2| < 3δ = ε} 被拆成四段；在目前所有假設（\code{0 < |x − 2| < δ}、\code{δ = ε/3}）下隨機取一組滿足假設的變數值，逐段驗：等式看兩邊差、不等式看方向。方向寫反的那一段標紅，並回傳一組帶數字的反例。

\begin{table}[h]
\centering\small\sffamily
\begin{tabularx}{\textwidth}{@{}lYl@{}}
\toprule
行 & 內容 & 判定 \\
\midrule
1 & 任取 $\varepsilon > 0$。 & \textcolor{moss}{綠} \\
2 & 取 $\delta = \varepsilon$。 & \textcolor{moss}{綠}（設定本身沒有錯） \\
3 & 假設 $0 < |x - 2| < \delta$。 & \textcolor{moss}{綠} \\
4 & 則 $|3x - 6| = 3|x - 2| < 3\delta = \varepsilon$。 & \textcolor{rust}{紅}：$3\delta = \varepsilon$ 不成立，反例 $\varepsilon = 0.3,\ \delta = 0.3$ \\
5 & 所以 $\lim_{x\to 2} 3x = 6$。 & \textcolor{gold}{黃}：結論沒有支撐 \\
\bottomrule
\end{tabularx}
\caption{一份 $\delta$ 取錯的 $\varepsilon$–$\delta$ 證明。紅在第 4 行，不是第 2 行——第 2 行的設定本身不是錯，錯的是它撐不起第 4 行。}
\end{table}

\section{推論引擎與骨架}
規則字典裡每一條定理有「需要的前提形狀」與「會吐出的結論形狀」。寫「由平均值定理，存在 $c \in (x, y)$ 使 $f(y) - f(x) = f'(c)(y - x)$」會被認出來；接著若跳過「因為 $f'(c) = 0$」直接寫「故 $f(x) = f(y)$」，結論不放行。每一題另有骨架檢查：一份 $\varepsilon$–$N$ 證明至少要有任取、取 $N$、假設 $n > N$、推導、結論五個角色，缺角色的證明就算每一行都綠也不算完成。

\section{怎麼知道它沒放水}
\file{validate\_proof\_lang.js} 對 43 題的參考證明做三種變異：刪一行、改一個符號、換一個方向。參考證明必須全綠；每一個變異版必須紅在正確的那一行——不是「有紅就好」。另外 68 條經典證明（步驟重排、填空決策點、自評）中有 8 條屬 Lean 層，其參考解在 CI 裡用 Lean 4 編譯，編不過就擋部署；其餘證明裡的定量宣稱（「這個積分等於 $\pi/4$」）由 \file{verify\_proof\_claims.js} 數值驗算。

\begin{decision}[標黃而不是標綠]
存在句、抽象函數、數論命題驗不了。最省事的做法是放行，最誠實的做法是標黃並寫「這一行看得懂但驗不了」。選了後者，代價是一份完全正確的證明可能不會全綠——說明頁直接告訴使用者這一點，而不是讓黃色看起來像錯。
\end{decision}

% ════════════════════════ 第 7 章 ════════════════════════
\chapter{引擎三：能力模型}
\section{要回答的問題}
別人給正確率；這裡要分得出「不會」「來不及」「壓力下垮掉」三種。同樣是 60\% 的正確率，一個人是每題都算到最後一秒才錯，另一個人是三秒就送出——處方完全不同。

\section{從標籤到技巧}
題庫有 324 個標籤，但那不是能力模型能用的座標系：標籤混了技巧、題型、題包來源、難度標記。\file{kernel/skill\_graph.js} 把標籤收斂成 80 個技巧節點並定義前置關係。一個技巧由「標籤命中且主題相符」決定——\code{taylor} 這個標籤同時出現在極限題與級數題，必須靠主題分開，否則極限題會被算成級數技巧。

\section{參數與理由}
\begin{table}[h]
\centering\small\sffamily
\begin{tabularx}{\textwidth}{@{}llY@{}}
\toprule
常數 & 值 & 為什麼 \\
\midrule
權重時間常數 & 30 天 & 30 天前的一題約值今天的 0.37 題 \\
先驗強度 & 6 & 等價於 6 次「不確定」的作答 \\
先驗正確率 & 0.45 & 偏低：沒證據就不給高分 \\
測準門檻 & 加權量 12 & 到這裡才算「測準了」；置信度低於 0.4 一律顯示「未測」而不是分數 \\
分流最低樣本 & 8 題 & 用原始題數，因為要能直接講給使用者聽 \\
看過解答才對 & 0.5 & 只給半分 \\
提示扣分 & 每層 0.15，下限 0.55 & 提示是拐杖，但不是零分 \\
「快」的門檻 & elapsed / timeLimit $<$ 0.6 & 相對題目的時限，不是絕對秒數 \\
「準」的門檻 & 正確率 $>$ 0.7 & \\
\bottomrule
\end{tabularx}
\caption{能力模型的常數。每一個都要能解釋，否則就會變成沒人敢動的魔術數字。}
\end{table}

\section{熟練度與置信度}
每筆作答的權重 $w = e^{-\Delta t / 30\text{ 天}}$，分數 $s \in [0, 1]$ 依對錯、是否看過解答、用了幾層提示調整。技巧的熟練度以先驗平滑：
\[
  \text{mastery} = \frac{\sum w_i s_i + 6 \times 0.45}{\sum w_i + 6},\qquad
  \text{confidence} = \min\Bigl(1, \frac{\sum w_i}{12}\Bigr)
\]
置信度不到 0.4 就顯示「未測」。這是產品層的決定：寧可留白，也不要給錯的能力數字。

\section{四個象限}
用「快不快」與「準不準」把每個技巧放進四格，每一格對應不同處方：
\begin{table}[h]
\centering\small\sffamily
\begin{tabular}{@{}llll@{}}
\toprule
象限 & 快 & 準 & 處方 \\
\midrule
反射區 & 是 & 是 & 拉高難度或進 Boss \\
會但慢 & 否 & 是 & 同技巧限時訓練，時限逐次縮短 \\
衝太快 & 是 & 否 & 練習模式 + 標註錯因 \\
還沒建立 & 否 & 否 & 看關鍵一句 + 三層提示 + 慢練 \\
\bottomrule
\end{tabular}
\caption{速度象限。「你會做，是壓力下垮掉」這句診斷來自限時與不限時正確率的差。}
\end{table}

\section{個人化的遺忘曲線}
30 天是預設的時間常數，但每個人忘得不一樣快。當一個人有 12 筆以上「隔多天回訪同一技巧」的實績，模型用回訪時的表現與預測的偏差校正半衰期倍率，範圍 0.6 到 1.8。首頁會說「依你 $n$ 次隔多天回訪的實績，你忘得比模型預設慢」。

\begin{decision}[回放，不存快照]
能力曲線的每一個歷史點都由紀錄與當時的時間重新算出，所以成長曲線不需要每天寫快照，使用者中斷三週再回來也不會斷線。代價是每次算要掃紀錄——第 11 章的分層儲存讓這件事在 5,000 筆逐題 tuple 上仍是毫秒級。
\end{decision}

% ════════════════════════ 第 8 章 ════════════════════════
\chapter{規劃器與抽題器}
\section{配方}
規劃器輸出的不是題目，是配方：每個角色佔多少比例，加上一句為什麼。三種基本配方：

\begin{table}[h]
\centering\small\sffamily
\begin{tabular}{@{}llll@{}}
\toprule
情境 & 角色與比例 & & \\
\midrule
維持手感 & 到期複習 0.6 & 一般練習 0.4 & \\
日常訓練 & 到期複習 0.35 & 弱點 0.4 & 新技巧 0.2 · 拉高難度 0.05 \\
新使用者 & 暖身 0.2 & 弱點 0.3 & 新技巧 0.3 · 難題 0.15 \\
\bottomrule
\end{tabular}
\caption{配方的角色比例。考前衝刺另有依考試日期與範圍調整的配方。}
\end{table}

\section{降級必須誠實}
一個新使用者沒有錯題可複習，一個只練過極限的人沒有級數的弱點資料，難度上限開 2 的人抽不到 R5。這些情況下配方一定有角色抽不滿——那時候必須降級補滿，不能留白。給使用者 12 題就是 12 題，少給 3 題比給錯 3 題更傷信任。降級路徑寫在 \code{meta.fallbacks}，畫面據此說「今天沒有到期的錯題，改多練兩題弱點」。

\begin{pitfall}[補位題掛在錯的名下]
第一版對全新使用者顯示「7 到期複習 · 8 弱點」——那是補位規則填出來的數字，對一個什麼都沒做過的人是假的。現在補位題不掛在任何角色名下，而且新手的第一份訓練固定 8 題、由淺入深。
\end{pitfall}

\section{新手保護與橋池}
剛從入門課程畢業的人直接丟進主線，第一題可能就是有理化或反三角——沒教過，信心當場塌掉。所以有一個「橋池」：校準後等級為 1、三個單元、不是文字判定題、而且沒有課裡沒教的標籤，131 題。畢業關從橋池抽 10 題、8 題過；畢業前與畢業後頭三局，任何訓練的池子都縮到橋池；頭四局不倒數。驗證器另外守橋池大小、課裡用的題全在橋池、排除的標籤真的存在於題庫（改名了就會安靜地漏進去）。

\section{每日一題與抽題去重}
每日一題由日期決定種子、全站同一題、一天一次；成績可以出成圖卡分享。所有抽題都經過等價側表去重，同一局不會出現只差記法的兩題。

% ════════════════════════ 第 9 章 ════════════════════════
\chapter{入門課程}
完全沒學過微積分的人打開站台會被 1,994 題淹死。入門課程補的是那一層：十六課，順序照一般課本——函數、極限與連續、微分、積分。

\begin{table}[h]
\centering\small\sffamily
\begin{tabularx}{\textwidth}{@{}lY@{}}
\toprule
單元 & 課 \\
\midrule
函數 & 函數是什麼（理）、常見的函數家族與合成（理） \\
極限與連續 & 極限是什麼、代進去變 0/0 怎麼辦、$x$ 跑到無窮遠、三個要背的極限、連續與中間值定理（理） \\
微分 & 導數就是斜率、微分公式、乘法律與除法律、連鎖律、$f'$ 的正負與極值（理）、均值定理（理） \\
積分 & 積分是微分的反操作、定積分是面積、黎曼和與基本定理（理） \\
\bottomrule
\end{tabularx}
\caption{十六課。標「理」的六課是理論課。}
\end{table}

兩種課：計算課是白話概念、一題來自題庫的逐步示範、小測（選錯會說為什麼）、三題不倒數的練習；理論課沒有題庫練習（題庫裡沒有、也不該有「定理說了什麼」這種反射題），示範是一段推導，小測至少三題，全對就算完成。逐步示範揭下一步之前會先讓人猜：真的下一步與同一題別的步驟二選一，猜錯照揭，只多一句「其實是這個」。每課完成後有一題同單元的 R2 延伸挑戰。

\file{validate\_course.js} 守：示範與練習必須是題庫裡的 R1 題、小測恰好一個正解且錯的選項都寫理由、課文裡「第 $n$ 課」的 $n$ 不超過總數、正解不能六成以上都放第一個、課文的每一條 TeX 都用本地 KaTeX 嚴格模式排一次。

% ════════════════════════ 第 10 章 ════════════════════════
\chapter{手寫計算紙}
\section{iPad 上的四個陷阱}
真正的使用情境是 iPad 加 Apple Pencil，那個情境有四個瀏覽器預設行為會直接毀掉它：長按選取（手掌碰到題目就整段反白、跳出拷貝選單）；手掌本身是 \code{touch} 事件，跟筆同時發生；Pencil 取樣 120–240 Hz 而 \code{pointermove} 一個 frame 只觸發一次，不撈 \code{getCoalescedEvents()} 就是多邊形；每次 move 重畫整塊板子是 $O(n^2)$，寫到第三行就看得出頓。

\section{繪製管線}
\begin{figure}[h]
\centering
\begin{tikzpicture}[
  node distance=5mm and 4mm,
  box/.style={draw=line,fill=paper,rounded corners=2pt,minimum height=10mm,text width=30mm,align=center,font=\sffamily\small},
  arr/.style={-{Latex[length=2mm]},draw=muted,line width=0.8pt}
]
\node[box] (a) {取樣\\{\scriptsize\color{muted}rawupdate / coalesced}};
\node[box,right=of a] (b) {去抖動\\{\scriptsize\color{muted}$< 0.6$ px 不收}};
\node[box,right=of b] (c) {一 frame 一次\\{\scriptsize\color{muted}rAF 合併}};
\node[box,right=of c] (d) {主畫布\\{\scriptsize\color{muted}只畫新增的曲線段}};
\node[box,below=of d] (e) {overlay\\{\scriptsize\color{muted}筆尖 + 預測點，每 frame 重畫}};
\draw[arr] (a) -- (b); \draw[arr] (b) -- (c); \draw[arr] (c) -- (d); \draw[arr] (c) |- (e);
\end{tikzpicture}
\caption{一筆從取樣到畫面。收筆時主畫布補到真正的筆尖，overlay 清空。}
\end{figure}

\begin{itemize}
  \item \textbf{平滑}：以相鄰兩個取樣點的中點為端點、取樣點本身為控制點畫二次曲線——這是把折線變成手寫感最便宜的做法，而且逐段可畫。第一段要從真正的起點出發，否則落筆的圓點跟線條之間有一小段空白。
  \item \textbf{增量}：用 \code{stroke.drawnTo} 記進度，只畫還沒畫過的那一段。
  \item \textbf{筆尖}：二次曲線要知道下一個取樣點才畫得出來，已提交的曲線永遠停在倒數第二個點。原本每個取樣都補一段直線到筆尖、下一個取樣再用曲線蓋掉——筆尖每秒抖兩百多次、繪製量三倍。現在筆尖那一段畫在 overlay，每 frame 清掉重畫，並接上瀏覽器的預測座標（Chrome 與 Android 有，Safari 沒有），墨水比筆尖早半步到。
  \item \textbf{寬度}：有壓感的筆用壓力，往回平均三點，落筆與收筆墊高（否則每一筆的尾巴收成針尖，一條分數線變成透鏡形）；沒有壓感的裝置用速度，作用範圍與壓感路徑相稱。
  \item \textbf{手掌}：看到筆就記時間，1.4 秒內的觸控一律忽略。
  \item \textbf{橡皮擦}：用 \code{destination-out} 真的挖掉墨水；拿背景色蓋過去會把方格紙的格線一起蓋掉。畫布保持透明，紙的方格與黑板底色都是 CSS。
\end{itemize}

\section{「有時候卡一下」}
\begin{pitfall}[真兇不是繪製]
使用者回報「有時候卡一下」。用 CDP Profiler 畫 40 筆、每筆 120 個取樣，主執行緒忙 1,564 ms，其中 \code{loadRecords} 151 ms——比繪製本身（\code{paintStrokeTail} 68 ms）還多。原因：每問一次筆寬或紙型就從 localStorage 讀整包紀錄再 JSON.parse，而繪製每 frame 都問。紀錄越大越卡。改成快取設定、存檔時清掉之後，\code{loadRecords} 從熱點消失，主執行緒忙 1,175 ms。
\end{pitfall}
教訓寫進規則：畫面上任何每 frame 跑的東西都不准呼叫 \code{loadRecords}；再有人說卡，先跑 \file{tools/profile\_board.js} 看長任務數與前十名自身時間，不要猜。

\section{測試怎麼測像素}
\file{e2e\_handwriting.js} 有 45 條斷言，全部打在畫布的像素與實際派發的 PointerEvent 上：手掌畫不出線、長按不跳選單、橡皮擦不蓋格線、復原不會整頁重繪、還在寫的時候墨水就已經到筆尖、每個取樣點只被畫一次（上線後的當機就是兩個監聽器都掛、每個取樣畫兩次）、寫過之後計算紙收得起來。合成事件不會自帶 coalesced 批次，測試手動補上——不補的話那條路徑在測試裡永遠不會發生，而那正是它漏到線上的原因。

% ════════════════════════ 第 11 章 ════════════════════════
\chapter{資料層與隱私}
\section{三層紀錄}
能力模型需要長時間的作答歷史，但完整的局明細上限只有 40 場——一個認真練的使用者兩週就洗掉了，等於每兩週失憶一次。解法是分三層，各自負責不同的時間尺度：

\begin{table}[h]
\centering\small\sffamily
\begin{tabularx}{\textwidth}{@{}llYl@{}}
\toprule
層 & 上限 & 用途 & 大小 \\
\midrule
history & 40 場 & 完整局與每題明細：結算頁、最近表現、明細檢視 & \\
attemptLog & 5,000 筆 & 極精簡逐題 tuple（不是物件）：能力模型、趨勢、百分位 & 約 250–300 KB，半年份 \\
sessions & 400 場 & 局級摘要：熱力圖、成長曲線 & 涵蓋一年以上 \\
\bottomrule
\end{tabularx}
\caption{Records v2。同一把 localStorage key、同一支正規化函式，只加頂層命名空間。}
\end{table}

舊使用者第一次載入時從 history 回填 attemptLog，不需要遷移腳本，也不會看到空白的趨勢圖。\file{validate\_records\_v2.js} 驗：history 截斷後能力模型仍保住 480 筆（只讀 history 只剩 96 筆）；兩台裝置的紀錄合併滿足交換律與冪等（80 + 60 → 140 筆，衝突以 \code{updatedAt} 新者為準）。

\section{手寫草稿與 IndexedDB}
筆畫陣列一題就能到幾十 KB，存 localStorage 會擠掉紀錄，所以草稿進 IndexedDB。存的理由只有一個：重做錯題時，「我上次是哪一步算錯的」這個問題的答案就在上次的草稿裡——上次答錯的草稿預設攤開對照，答對的收起（攤開就變照抄）。

\section{隱私承諾與它的測試}
政策上寫的每一條在 \file{validate\_privacy.js} 裡都有靜態檢查：
\begin{itemize}
  \item 關掉分析不只是不送事件，\textbf{連 gtag.js 都不載}——光是載入那支 script 就已經對 Google 發出請求；檢查器確認 \code{analyticsEnabled()} 的判斷在載入 script 之前。
  \item 清除資料掃的是整個 \code{buzzcalculus.} 命名空間，不是一份固定清單——自動備份的 key 帶時間戳，固定清單永遠列不完；E2E 曾抓到「按了清除，一份完整備份還躺在 localStorage 裡」。
  \item 錯誤遙測只送型別、檔名、行號，不送訊息與堆疊——JS 引擎組出來的訊息常把變數的值印進去，那可能是題目或使用者的答案。
  \item 回報題目不能只寫進使用者自己的瀏覽器：那比沒有回報按鈕更糟，它製造了「已經反映過」的錯覺。
\end{itemize}

\section{零後端的社群功能}
對戰碼 \code{BZD1:417.882.…\textasciitilde 8.225} 裝的是題目的永久編號與發戰帖者的成績，貼給朋友、對方貼碼應戰、結算頁直接分勝負，不用帳號不用網路。成就分享卡與每日一題成績卡在本機 canvas 畫成 PNG，手機走系統分享（\code{navigator.share} 帶檔案），桌機下載——資料一律不出裝置。

% ════════════════════════ 第 12 章 ════════════════════════
\chapter{驗證文化與 CI}
\section{四十一支驗證器}
\begin{longtable}{@{}p{4.6cm}p{10.6cm}@{}}
\caption{驗證器與它守的東西。全部只用 Node 內建模組。}\\
\toprule
驗證器 & 守什麼 \\
\midrule
\endfirsthead
\toprule
驗證器 & 守什麼 \\
\midrule
\endhead
verify\_answers & 從題幹獨立驗算標準答案，不符必須為零 \\
validate\_problems & 題目欄位完整、答案格式合法、時限與難度存在 \\
validate\_answer\_checker & 判分器對固定的正反例判得對 \\
validate\_content\_metadata & 永久編號、三軸、來源聲明齊全；名校字眼不准標原創 \\
validate\_calculus\_only & 題庫必須是純微積分 \\
validate\_katex & 每一題的題幹真的拿 KaTeX 渲染一遍 \\
validate\_hints / derived\_hints / graph\_choices / worksheets & 提示品質、推導式提示重新驗證、選圖題與作圖表的把關 \\
validate\_daily\_one & 每日一題釘選表；歷史可回放 \\
validate\_training\_packs / path\_nodes / problem\_selection & 題包非空、主線每一關抽得出題、抽題規則 \\
validate\_skill\_graph & 80 個技巧節點與前置關係一致 \\
validate\_ability\_model & 能力模型在合成紀錄上的行為（衰減、留白、象限） \\
validate\_planner & 配方帶 why、降級路徑、等價去重 \\
validate\_records\_v2 & 分層截斷保住資料、合併交換律與冪等 \\
validate\_proof\_lang & 參考證明全綠、三種變異紅在正確的行 \\
verify\_proof\_claims & 經典證明裡的定量宣稱數值驗算 \\
validate\_course & 入門課程的示範、練習、小測、TeX、橋池 \\
validate\_golden & 黃金檔：判分、繪圖指令序列、驗算方法的中間結果 \\
validate\_analytics & 事件表與程式碼一致、不上報使用者輸入 \\
validate\_privacy & 政策裡的每一條承諾 \\
validate\_offline\_assets / app\_shell / version & 快取清單、出貨頁面、版本號 \\
validate\_performance\_budget & app.js、kernel、樣式、題庫、script 數各有預算 \\
validate\_public\_claims & README 的 21 個數字從程式碼算出來 \\
spec\_snapshot & 規格文件的現況快照必須與程式碼一致 \\
smoke\_app\_render & 在 Node 裡 render 一次首頁 \\
regen\_all --check & 產生器的產物與版控同步 \\
\bottomrule
\end{longtable}

\section{九支端對端測試}
E2E 不用 puppeteer：Node 22 內建 WebSocket 與 fetch，而 Chrome DevTools Protocol 就是「連上 ws、送 JSON」，\file{tools/lib/cdp.js} 兩百多行就夠。斷言原則：只斷言使用者看得見、點得到的東西——文字要有面積且 opacity 不是 0，互動要真的觸發事件委派，console 不能有錯、資源不能有 404。其他驗證器比對的是 render 出來的 HTML 字串，那一層看不到「畫面是白的」——實際抓到過結算頁 DOM 有兩萬字、畫面零字。

\begin{table}[h]
\centering\small\sffamily
\begin{tabular}{@{}lrl@{}}
\toprule
套件 & 斷言 & 走的路 \\
\midrule
main\_flow & 34 & 開局、答題、結算、錯題進本、四個分頁、快捷鍵、清除資料 \\
storage & 16 & localStorage 額度爆掉時的救援、匯出、離線重載 \\
handwriting & 45 & 畫布像素與 PointerEvent（第 10 章） \\
mobile & 41 & 390×844 真視窗：橫向溢出、觸控目標 44 px、iPad 鍵盤 \\
layout & 35 & 外殼幾何：分頁列貼底、側欄不蓋內容 \\
responsive & 93 態 & 9 種尺寸 × 9 頁 \\
proof\_lang & 37 & 教學到題目、打字即時三色、提交紀錄 \\
beginner & 72 & 定位、新手保護、十六課、畢業關、橋池 \\
graph\_interactive & 10 & SVG 上的點位與拖曳 \\
\bottomrule
\end{tabular}
\caption{端對端測試。本機用 \file{e2e\_all.js} 並行跑，九支約十分鐘。}
\end{table}

\section{CI 上看不見的失敗}
GitHub Actions 的 job log 要認證才抓得到，而 check-run 的 annotations 走公開 API 就讀得到。所以每一條失敗的斷言都印成 \code{::error title=…::…}，30 支驗證器由一支執行器逐支跑、紅的那支變成 annotation。「CI 紅、本機綠」時，這是唯一能看到 runner 上哪條倒了的方法。

\begin{pitfall}[CI 從 9 月 1 日紅到 9 月 16 日]
部署與測試在同一個 job，任何一支 E2E 紅站台就不更新，兩週的工作使用者一個都沒看到。五支 E2E 全是同一種病：\code{sleep(500)} 這種「開發機剛好夠」的固定等待，在慢好幾倍的 runner 上輸給 debounce、rAF render 與 950 ms 的自動前進。修法一律是「輪詢等元素真的出現」，斷言本身不改；另外加了 CPU 節流開關在本機重現 runner 的節奏。還有兩個假陽性：離線重載其實只是改 hash 沒有真的重載；結算頁的彩帶粒子飛完停在 opacity 0，被當成「內容被藏住」。
\end{pitfall}

\section{公開的數字必須是真的}
README 裡的題數、驗算題數、技巧數用標記包起來，\file{validate\_public\_claims.js} 從程式碼算出來比對；規格文件開頭的現況快照（題數、app.js 行數、驗證器數）由 \file{spec\_snapshot.js} 產生並在 CI 檢查一致。數字錯掉的話，讀的人會拿它當事實去做取捨。

% ════════════════════════ 第 13 章 ════════════════════════
\chapter{效能預算與模組化}
\section{預算}
\begin{table}[h]
\centering\small\sffamily
\begin{tabular}{@{}lrrr@{}}
\toprule
項目 & 現況 & 預算 & 使用率 \\
\midrule
app.js 主程式 & 722 KB & 720 KB → 搬模組後回到預算內 & — \\
kernel 邏輯模組 & 271 KB & 280 KB & 97\% \\
題庫資料 & 995 KB & 2,400 KB & 41\% \\
styles.css & 305 KB & 310 KB & 98\% \\
其他 src 腳本 & 172 KB & 300 KB & 57\% \\
script 標籤 & 65 & 65 & 100\% \\
\bottomrule
\end{tabular}
\caption{效能預算。撞頂時不調數字，搬整段出去；要調就寫清楚為什麼。}
\end{table}

\section{撞頂就搬，不調數字}
這一輪兩次撞到 app.js 的 720 KB：第一次把週報與成就分享卡（230 行 canvas 程式）搬成 \file{share\_cards.js}；第二次把錯因處方、段考預測、模式推薦三支純函式搬進 \file{kernel/planner.js}，app 端的依賴用 \code{deps} 注入。兩個坑：搬出去的東西若被 app.js 的 hooks 物件用簡寫引用（\code{weeklyShareData,}），文字搜尋抓不到，要靠型別檢查；模組在 IIFE 裡的 \code{create} 呼叫執行時，其後才宣告的 \code{const} 還在暫時死區，站台 URL 這種常數要用 getter 傳。

% ════════════════════════ 第 14 章 ════════════════════════
\chapter{結語}
這個專案裡真正難的東西，都是「在沒有後端的限制下把一件本來要靠伺服器的事做對」：判分、證明、能力模型、手寫。但讓它能被信任的不是任何一個引擎，而是把每一條承諾綁進 CI 的習慣——答案錯一題會紅、政策說謊會紅、README 的數字錯會紅、畫面白掉會紅。

還沒做的事：QR code 搬紀錄（紀錄幾十 KB 塞不進 QR，要先壓縮加分片）；app.js 拆首屏（目前的策略是撞預算就搬整段，總有一天要正面處理）；白話證明對存在句與抽象函數的三分之一；能力模型的遺忘曲線需要更多回訪實績才能個人化。這些都寫在規格文件裡，不是空話。

\appendix
\chapter{常數表}
\begin{table}[h]
\centering\small\sffamily
\begin{tabularx}{\textwidth}{@{}p{3.4cm}YY@{}}
\toprule
常數 & 意義 & 值 \\
\midrule
取樣點 & 判分引擎的單變數取樣 & 0.3137, 0.7211, 1.2345, 1.9871, 3.3013, −0.6180, −1.3247, −2.1069 \\
多變數取樣 & 線性同餘產生 & 24 點，種子 20260914 \\
判分容差 & 相對 & $\max(10^{-6}, 10^{-5}|a|)$ \\
判得穩 & 至少幾點有定義 & 3 \\
權重時間常數 & 能力模型 & 30 天 \\
先驗 & 強度 / 正確率 & 6 / 0.45 \\
測準門檻 & 加權作答量 & 12 \\
顯示門檻 & 置信度 & 0.4 \\
提示扣分 & 每層 / 下限 & 0.15 / 0.55 \\
遺忘倍率 & 範圍 / 最少樣本 & 0.6–1.8 / 12 \\
history / attemptLog / sessions & 上限 & 40 / 5,000 / 400 \\
橋池 & 校準後 R1、三單元、非文字題、無未教標籤 & 131 題 \\
畢業關 & 題數 / 過關 & 10 / 8 \\
新手保護 & 不倒數的局數（沒學過 / 高中先修 / 大一） & 4 / 3 / 1 \\
手掌視窗 & 看到筆之後忽略觸控 & 1,400 ms \\
去抖動 & 相鄰取樣距離 & 0.6 px \\
壓感平均 & 往回幾點 & 3 \\
自動前進 & 答對之後 & 950 ms \\
CI 等待倍率 & 固定 sleep 乘 & 3 \\
\bottomrule
\end{tabularx}
\end{table}

\chapter{詞彙表}
\begin{description}[style=nextline,leftmargin=0pt,itemsep=4pt]
  \item[R1–R6] 難度等級，由三軸（步數、冷門度、負荷）算出。R1 暖身、R2 基礎、R3 標準、R4 進階、R5 Boss、R6 Boss+。
  \item[技巧（skill）] 能力模型的座標，80 個，由 324 個標籤收斂而成，有前置關係。
  \item[配方（recipe）] 規劃器的輸出：每個角色佔多少題，加一句為什麼。
  \item[降級（fallback）] 抽題器抽不滿某個角色時的補滿路徑，記在 \code{meta.fallbacks}。
  \item[橋池] 入門課程與主線之間的題池：只出課裡教過的 R1 題。
  \item[三色] 白話證明的逐行判定：綠（驗過成立）、紅（驗過不成立）、黃（看得懂但驗不了）。
  \item[黃金測試] 錄下固定輸入的完整輸出序列，任何改動都要比對。
  \item[側表] 由產生器從題庫算出、進版控的資料表：永久編號、三軸、來源、驗算結果、等價組。
  \item[annotation] GitHub Actions 的 \code{::error} 標記，不用認證就能從公開 API 讀到。
\end{description}

\end{document}
